\textbf{Pour cette première partie, aucune justification n'est demandée et une seule réponse est possible par question.} Les justifications ci-dessous sont fournies à titre pédagogique.

Une réponse fausse ou l'absence de réponse n'enlève aucun point : il ne faut donc jamais laisser une question sans réponse.

\subsection*{Question 1 \quad Réponse B}

On range d'abord les 7 notes dans l'ordre croissant :
\[8 \quad 9 \quad 9 \quad \mathbf{11} \quad 12 \quad 15 \quad 17.\]
L'effectif est impair (7 notes) : la médiane est la valeur centrale, c'est-à-dire la
4\up{e} note, soit 11.

\subsection*{Question 2 \quad Réponse D}

La droite coupe l'axe des ordonnées en 6 : l'ordonnée à l'origine est $b = 6$. Elle passe
aussi par le point de coordonnées $(2 ; 0)$, donc son coefficient directeur vaut
\[a = \dfrac{0 - 6}{2 - 0} = -3.\]
D'où $f(x) = -3x + 6$.

\subsection*{Question 3 \quad Réponse B}

On reconnaît l'identité remarquable $a^2 - b^2 = (a-b)(a+b)$ avec $a = x$ et $b = 3$ :
\[(x-3)(x+3) = x^2 + 3x - 3x - 9 = x^2 - 9.\]

\subsection*{Question 4 \quad Réponse A}

Prendre 30\,\% d'une quantité revient à la multiplier par $\dfrac{30}{100} = 0,3$ :
\[0,3 \times 12 = 3,6.\]

\subsection*{Question 5 \quad Réponse C}

Diviser par une fraction revient à multiplier par son inverse :
\[A = \dfrac{3}{7} \div \dfrac{4}{5} = \dfrac{3}{7} \times \dfrac{5}{4} = \dfrac{15}{28}.\]

\subsection*{Question 6 \quad Réponse C}

Il s'agit d'une situation de proportionnalité : 1~mètre correspond à 1,1~yard, donc
\[300 \times 1,1 = 330.\]
300~mètres représentent environ 330~yards.

\subsection*{Question 7 \quad Réponse D}

On écrit $2 = 2^1$ et on applique $a^m \times a^n = a^{m+n}$ :
\[2 \times 2^5 = 2^1 \times 2^5 = 2^{1+5} = 2^6.\]

\subsection*{Question 8 \quad Réponse B}

On remplace $T_C$ par $-10$ dans la formule, en respectant la priorité de la
multiplication :
\[T_F = 1,8 \times (-10) + 32 = -18 + 32 = 14.\]
Une température de $-10$~\textdegree C correspond à 14~\textdegree F.
